\documentclass[12pt,a4paper]{article}
\usepackage[utf8]{inputenc}
\usepackage[T1]{fontenc}
\usepackage{lmodern}
\usepackage{geometry}
\usepackage{graphicx}
\usepackage{float}
\usepackage{amsmath}
\usepackage{array}
\usepackage{booktabs}
\usepackage{caption}
\usepackage{hyperref}
\usepackage{placeins}

\geometry{left=2.5cm,right=2.5cm,top=2.5cm,bottom=2.5cm}
\hypersetup{colorlinks=true, linkcolor=black, urlcolor=blue, citecolor=black}
\renewcommand{\contentsname}{Sadržaj}
\renewcommand{\listfigurename}{Popis slika}
\renewcommand{\listtablename}{Popis tabela}
\renewcommand{\figurename}{Slika}
\renewcommand{\tablename}{Tabela}
\setlength{\parindent}{0pt}
\setlength{\parskip}{0.75em}
\sloppy

\begin{document}

\begin{titlepage}
\centering
\vspace*{1.8cm}
{\Large UNIVERZITET U SARAJEVU\par}
\vspace{0.7cm}
{\large ELEKTROTEHNIČKI FAKULTET\par}
{\large ODSJEK ZA AUTOMATIKU I ELEKTRONIKU\par}
\vspace{4.0cm}
\rule{0.82\textwidth}{0.4pt}
\vspace{0.8cm}

{\LARGE\bfseries Modeliranje i integracija solarne elektrane sa inverterom u dTwin okruženju\par}
\vspace{0.35cm}
{\large Predmet: Sistemi i režimi rada elektroenergetskih sistema\par}
\vspace{0.8cm}
\rule{0.82\textwidth}{0.4pt}
\vspace{1.2cm}

{\large Student: Hana Zaimović\par}
{\large Broj indeksa: 19658\par}
\vfill
{\large Sarajevo, juli 2026.\par}
\end{titlepage}

\pagenumbering{roman}

\section*{Sažetak}
U ovom radu prikazana je implementacija C++/dTwin plugin-a za modeliranje solarne elektrane sa inverterom i njeno priključenje na odabrani čvor elektroenergetskog sistema. Aplikacija omogućava izbor testnog elektroenergetskog sistema, izbor priključnog čvora, učitavanje XML konfiguracije, zadavanje scenarija rada, proračun aktivne i reaktivne snage invertera, formiranje matričnih podataka i generisanje dTwin modela.

Korišteni testni sistemi su case9, case30, case118 i case300. Za odabrani slučaj korisnik može ručno odabrati ili automatski preporučiti čvor priključenja PV invertera. Model invertera predstavljen je kao dinamički DAE model sa vremenskom konstantom, zbog čega se nakon numeričkog rješavanja u dTwin-u dobijaju vremenski odzivi aktivne AC snage, reaktivne snage i dostupne DC snage. Dodatno je omogućen prikaz rezultata u stvarnim jedinicama MW/MVAr i u per-unit sistemu.

\textbf{Ključne riječi:} solarna elektrana, PV inverter, dTwin, natID matrice, DAE model, per-unit sistem

\section*{Abstract}
This report presents the implementation of a C++/dTwin plugin for modelling a photovoltaic power plant with a grid-tied inverter and connecting it to a selected bus of a power system test case. The application supports XML-based configuration, test case selection, manual or scenario-based operating point definition, calculation of active and reactive inverter power, matrix information preview and automatic generation of a dTwin model.

The implemented inverter model is represented as a dynamic DAE model with a first-order response. After numerical simulation in dTwin, the model produces time-domain responses of AC active power, reactive power and available DC power. The results can be shown both in engineering units and in the per-unit system. The implementation is tested on case9, case30, case118 and case300 networks and uses natID dense and sparse matrix structures.

\textbf{Keywords:} photovoltaic power plant, PV inverter, dTwin, natID matrices, DAE model, per-unit system

\clearpage
\tableofcontents
\clearpage
\listoffigures
\clearpage
\listoftables
\clearpage
\pagenumbering{arabic}

\section{Uvod}
Savremeni elektroenergetski sistemi sve više uključuju distribuirane izvore energije, posebno solarne elektrane povezane na mrežu preko energetskih pretvarača. Za razliku od klasičnih sinhronih generatora, solarna elektrana primarno proizvodi istosmjernu snagu, dok je za priključenje na izmjeničnu mrežu potreban inverter. Zbog toga se solarna elektrana u mrežnim proračunima posmatra kao kombinacija PV polja, upravljačkog sistema i invertera koji definiše aktivnu i reaktivnu snagu na mjestu priključenja.

Cilj projekta je razvoj softverskog alata koji iz XML konfiguracije solarne elektrane i XML opisa testne mreže generiše dTwin model. Projekat obuhvata C++ implementaciju, CMake konfiguraciju, natID dense i sparse matrice, GUI plugin za dTwin i testiranje na više mreža različite veličine.

\section{Opis problema}
Solarna elektrana se u ovom radu posmatra kao izvor koji proizvodi dostupnu DC snagu zavisnu od iradijanse i temperature. Ta snaga se preko invertera pretvara u AC snagu koja se injektuje u elektroenergetsku mrežu. Korisnik bira čvor mreže na koji će se inverter priključiti, a program zatim provjerava izabrani čvor, računa radnu tačku invertera i generiše dTwin model.

Problem se može opisati kroz sljedeće korake:
\begin{enumerate}
\item učitati konfiguraciju PV invertera iz XML fajla,
\item učitati odabrani testni mrežni slučaj,
\item izabrati priključni čvor,
\item izračunati dostupnu DC snagu, AC snagu i reaktivnu snagu,
\item formirati podatke za dense i sparse matrice,
\item generisati dTwin model,
\item riješiti model u dTwin-u i prikazati odzive kroz vrijeme.
\end{enumerate}

U prikazanom primjeru korišten je case30, sa odabranim busom 8, iradijansom $G=400\,\mathrm{W/m^2}$, temperaturom $T=20^\circ\mathrm{C}$ i faktorom snage $\mathrm{PF}=0.95$.

\section{Teorijska osnova}
\subsection{Kompleksna snaga}
U izmjeničnim elektroenergetskim sistemima aktivna i reaktivna snaga se najčešće posmatraju zajedno kroz kompleksnu snagu:
\[
S = P + jQ.
\]
Aktivna snaga $P$ predstavlja korisnu energiju koja se pretvara u rad, toplotu ili drugi oblik energije, dok reaktivna snaga $Q$ opisuje razmjenu energije između električnog i magnetnog polja. Za inverter u ovom projektu aktivna snaga $P_{ac}$ predstavlja AC snagu koju inverter injektuje u mrežu, dok $Q_{ac}$ predstavlja reaktivnu snagu određenu faktorom snage.

U implementaciji je korišten pojednostavljeni model:
\[
P_{ac}=\min(P_{dc},P_{ac,rated}),
\qquad
Q_{ac}=P_{ac}\tan(\arccos(\mathrm{PF})).
\]

\subsection{Solarni izvor i DC snaga}
Solarna elektrana proizvodi DC snagu koja zavisi od uslova okoline. Najvažniji parametri su iradijansa i temperatura ćelije. U implementaciji je korišten model:
\[
P_{dc}=P_{rated}\frac{G}{G_{ref}}\left(1+k_T(T_{cell}-T_{ref})\right),
\]
gdje su $P_{rated}$ nazivna DC snaga PV polja, $G$ trenutna iradijansa, $G_{ref}$ referentna iradijansa, $T_{cell}$ temperatura ćelije, $T_{ref}$ referentna temperatura i $k_T$ temperaturni koeficijent.

\subsection{Dinamički model invertera}
Za potrebe dTwin modela inverter nije modeliran kao prekidački PWM uređaj, nego kao dinamički izvor snage sa prvorednim odzivom. Time se dobija glatka promjena snage kroz vrijeme umjesto trenutne promjene iz nule u konačnu vrijednost:
\[
\frac{dp_{ac}}{dt}=\frac{P_{ac,ref}-p_{ac}}{\tau},\qquad
\frac{dq_{ac}}{dt}=\frac{Q_{ac,ref}-q_{ac}}{\tau},\qquad
\frac{dp_{dc}}{dt}=\frac{P_{dc,ref}-p_{dc}}{\tau}.
\]
U modelu je korištena vremenska konstanta $\tau=0.12\,\mathrm{s}$.

\begin{figure}[H]
\centering
\includegraphics[width=0.72\textwidth]{figures/fig03_dtwin_model_code.png}
\caption{Generisani dTwin DAE model za PV inverter sa parametrima i diferencijalnim jednačinama.}
\end{figure}

\subsection{Per-unit sistem}
Per-unit sistem se koristi da bi se električne veličine prikazale u odnosu na bazne vrijednosti:
\[
x_{pu}=\frac{x}{x_{base}}.
\]
U ovom projektu p.u. prikaz se koristi kod dodatnog dTwin grafa. Aktivna AC snaga i reaktivna snaga normalizuju se prema nazivnoj AC snazi invertera, dok se DC snaga normalizuje prema nazivnoj DC snazi fotonaponskog polja:
\[
p_{ac,pu}=\frac{p_{ac}}{P_{ac,n}},\qquad
q_{ac,pu}=\frac{q_{ac}}{P_{ac,n}},\qquad
p_{dc,pu}=\frac{p_{dc}}{P_{dc,n}}.
\]

\subsection{Ybus matrica i natID strukture}
Elektroenergetska mreža se može opisati metodom napona čvorova:
\[
YV=I,
\]
gdje je $Y$ sistemska matrica admitansi, $V$ vektor napona čvorova, a $I$ vektor injektovanih struja. Budući da elektroenergetske mreže obično nisu potpuno povezane, Ybus matrica ima veliki broj nultih elemenata, pa je prirodno koristiti sparse reprezentaciju. U projektu se zato prikazuju i koriste natID dense i sparse matrične strukture.

\section{Implementacija GUI-ja i plugin-a}
Razvijeni GUI nosi naziv \textit{PV inverter converter}. Interfejs je organizovan tako da korisnik prvo definiše ulazne fajlove i scenario, zatim pregleda mrežu i rezultate, a na kraju generiše dTwin model.

\begin{figure}[H]
\centering
\includegraphics[width=0.55\textwidth]{figures/fig07_plugin_import_menu.png}
\caption{dTwin import meni sa registrovanim pluginom PV inverter converter.}
\end{figure}

Polje \textit{Config XML} služi za izbor konfiguracijskog fajla solarne elektrane i invertera. Polje \textit{Case XML} služi za izbor testne mreže, a polje \textit{Output} određuje gdje će biti generisan izlazni dTwin model. Dugmad case9, case30, case118 i case300 omogućavaju brzo učitavanje testnih slučajeva.

\begin{figure}[H]
\centering
\includegraphics[width=0.95\textwidth]{figures/fig01_gui_initial.png}
\caption{Početni izgled GUI-ja prije učitavanja konfiguracije i izvršavanja preview proračuna.}
\end{figure}

Dugmad Sunny, Cloudy, Hot i Low PF predstavljaju unaprijed definisane scenarije rada. Manual dio omogućava ručnu promjenu priključnog busa, iradijanse, temperature i faktora snage. Dugme \textit{Recommend bus} automatski bira demonstracioni PQ/load bus pogodan za priključenje PV invertera.

Dugme \textit{Preview} izvršava proračun bez generisanja finalnog dTwin modela. Nakon toga GUI popunjava informacije o učitanom case-u, odabranom busu, PV/inverter parametrima, rezultatima i matričnim podacima. Dugme \textit{Convert} generiše dTwin model u stvarnim jedinicama, a \textit{Convert + p.u.} dodaje i per-unit prikaz.

\begin{figure}[H]
\centering
\includegraphics[width=0.95\textwidth]{figures/fig02_gui_preview_case30.png}
\caption{GUI nakon preview proračuna za case30, priključni bus 8 i Cloudy scenario.}
\end{figure}

\section{Rezultati}
Za prikaz rezultata korišten je case30, priključni bus 8 i Cloudy scenario. U tom scenariju je $G=400\,\mathrm{W/m^2}$, $T=20^\circ\mathrm{C}$ i $\mathrm{PF}=0.95$. Izračunate vrijednosti su $P_{dc}=3.264\,\mathrm{MW}$, $P_{ac}=3.264\,\mathrm{MW}$ i $Q_{ac}=1.0728\,\mathrm{MVAr}$.

\begin{table}[H]
\centering
\caption{Rezultati testnih slučajeva za Cloudy scenario}
\begin{tabular}{cccccc}
\toprule
Case & Čvorovi & $P_{ac}$ & $Q_{ac}$ & Dense & Sparse nonzeros \\
\midrule
case9 & 9 & 3.264 MW & 1.0728 MVAr & 9x9 & 26 \\
case30 & 30 & 3.264 MW & 1.0728 MVAr & 30x30 & 89 \\
case118 & 118 & 3.264 MW & 1.0728 MVAr & 118x118 & 353 \\
case300 & 300 & 3.264 MW & 1.0728 MVAr & 300x300 & 899 \\
\bottomrule
\end{tabular}
\end{table}

\begin{figure}[H]
\centering
\includegraphics[width=0.90\textwidth]{figures/fig04_dtwin_power_plot.png}
\caption{Vremenski odziv aktivne AC snage, reaktivne snage i dostupne DC snage u dTwin-u.}
\end{figure}

\begin{figure}[H]
\centering
\includegraphics[width=0.90\textwidth]{figures/fig05_dtwin_pu_plot.png}
\caption{Normalizovani per-unit prikaz odziva snaga PV invertera.}
\end{figure}

Rezultati pokazuju da se isti model invertera može povezati na različite mreže bez promjene osnovne konfiguracije solarne elektrane. Razlika između testnih slučajeva prvenstveno se vidi kroz dimenzije matrica i broj nenultih elemenata u sparse matrici. Dinamički odziv snaga ima oblik prvog reda, pri čemu se sve posmatrane snage približavaju izračunatim referentnim vrijednostima.

Vrijednosti napona prije i poslije priključenja PV invertera prikazane u GUI-ju predstavljaju indikativnu procjenu uticaja invertera na izabrani čvor. One nisu rezultat kompletnog Newton-Raphson proračuna tokova snaga, jer ulazni XML testni slučajevi u ovoj verziji projekta sadrže osnovne podatke o mreži, ali ne i kompletne bus, branch i generator podatke potrebne za puni AC power flow.

\section{Zaključak}
U okviru projekta realizovan je C++ konverter za model solarne elektrane sa inverterom. Model koristi XML konfiguraciju, podržava standardne testne slučajeve elektroenergetskih mreža, formira natID dense i sparse matrice i generiše dTwin model za dinamičku simulaciju. Plugin GUI omogućava praktičan izbor ulaznih fajlova, scenario pristup, ručnu izmjenu parametara, preview rezultata i generisanje dTwin prikaza u apsolutnim i per-unit vrijednostima.

Implementirani model predstavlja osnovu za dalja proširenja, kao što su detaljniji model regulatora invertera, napredniji proračun uticaja priključnog čvora na naponski profil i direktna integracija sa kompletnim proračunom tokova snaga.

\clearpage
\begin{thebibliography}{9}
\bibitem{milano} F. Milano, \textit{Power System Modelling and Scripting}, Springer, 2010.
\bibitem{natid} natID SDK dokumentacija i primjeri, idzafic/natID.
\bibitem{dtwin} dTwin dokumentacija i primjeri plugin integracije, idzafic/dTwin.
\bibitem{ieee} IEEE test systems: case9, case30, case118 i case300.
\end{thebibliography}

\end{document}
